\documentclass[11pt,a4paper]{article}

% -----------------------------------------------------------------------
% Packages
% -----------------------------------------------------------------------
\usepackage{amsmath,amsthm,amssymb,mathtools}
\usepackage{geometry}
\usepackage[T1]{fontenc}
\usepackage{lmodern}
\usepackage{microtype}
\usepackage{booktabs}
\usepackage{enumitem}
\usepackage{hyperref}
\usepackage{url}

\geometry{margin=1in}

% -----------------------------------------------------------------------
% Theorem environments
% -----------------------------------------------------------------------
\theoremstyle{plain}
\newtheorem{theorem}{Theorem}
\newtheorem{proposition}[theorem]{Proposition}
\newtheorem{lemma}[theorem]{Lemma}
\newtheorem{corollary}[theorem]{Corollary}

\theoremstyle{definition}
\newtheorem{definition}[theorem]{Definition}
\newtheorem{example}[theorem]{Example}

\theoremstyle{remark}
\newtheorem{remark}[theorem]{Remark}

% -----------------------------------------------------------------------
% Macros
% -----------------------------------------------------------------------
\newcommand{\floor}[1]{\left\lfloor #1 \right\rfloor}
\newcommand{\val}{\nu}
\newcommand{\ones}{N_1}
\newcommand{\zeros}{N_0}

\title{Asymptotic and structural results for OEIS sequence A112509}
\author{Ben Clare}
\date{April 2026}

\begin{document}
\maketitle

\begin{abstract}
For a binary string $s$ of length $n$, let $f(s)$ denote the number of
distinct non-negative integers whose binary representations appear as
contiguous substrings of $s$ (leading zeros suppressed), and set
$a(n)=\max_{|s|=n}f(s)$ (OEIS A112509~\cite{oeis}).
We prove three results.
First, $\lim_{n\to\infty}a(n)/n^2=\frac{1}{2}$; more precisely,
for all $n\ge64$,
\[
  \tfrac{1}{2}n^2 - 15\,n^{3/2} \;\le\; a(n) \;\le\; \tfrac{n(n+1)}{2},
\]
so $a(n)/n^2 = \frac{1}{2}+O(n^{-1/2})$.
Second, every optimal string has ones-density $1-O(n^{-1/4})$;
explicit near-optimal strings achieving this bound are constructed.
Third, $a(n)$ is convex: $a(n+1)-2a(n)+a(n-1)\ge0$ for all $n\ge2$,
so the first differences $\Delta a(n)$ are nondecreasing.
\end{abstract}

\section{Introduction}

\subsection{The sequence}

For $s\in\{0,1\}^n$, let $f(s)$ denote the number of distinct
non-negative integers whose binary representations appear as contiguous
substrings of $s$ (leading zeros suppressed: $\mathtt{01}=1$,
$\mathtt{00}=0$).  Set $a(n)=\max_{|s|=n}f(s)$.  This is OEIS
A112509~\cite{oeis}, beginning
\[
  1,3,5,7,10,13,17,22,27,33,40,47,55,64,73,83,94,106,118,131,\ldots
\]

\subsection{Numerical evidence}

Exact values for $n\le80$ are due to M.\ Fuller and appear in the OEIS
b-file~\cite{oeis}; no algorithm or proof of completeness was published.
Heuristic values for $81\le n\le150$ were computed in~\cite{github}
using multiple independent Metropolis--Hastings and template-search
methods, with consistent agreement between methods but no formal proof
of optimality.
Certified lower bounds for much larger $n$ were obtained via
suffix-array-guided local search~\cite{github}.
Table~\ref{tab:evidence} records selected values.

\begin{table}[ht]
\centering
\begin{tabular}{rrl}
\toprule
$n$ & $a(n)$ or lower bound & $a(n)/n^2\ge$ \\
\midrule
$50$            & $899$                                    & $0.3596$ \\
$100$           & $3{,}875$                                & $0.3875$ \\
$150$           & $9{,}070$                                & $0.4031$ \\
$10^6$          & $498{,}396{,}662{,}200$                  & $0.4984$ \\
$10^9$          & $499{,}943{,}911{,}277{,}037{,}650$      & $0.49994$ \\
$2\times10^9$   & $1{,}999{,}841{,}302{,}297{,}432{,}600$ & $0.49996$ \\
\bottomrule
\end{tabular}
\caption{Selected values and lower bounds for $a(n)$.
  Values for $n\le80$ are from Fuller~\cite{oeis};
  values for $81\le n\le150$ are heuristic computational estimates~\cite{github};
  large-$n$ entries are certified lower bounds~\cite{github}.}
\label{tab:evidence}
\end{table}

\subsection{Main results}

\begin{theorem}[Limit]\label{thm:limit}
$\displaystyle\lim_{n\to\infty}\dfrac{a(n)}{n^2}=\dfrac{1}{2}$.
\end{theorem}

\begin{theorem}[Quantitative lower bound]\label{thm:quant}
For every $n\ge64$,
\[
  a(n)\;\ge\;\frac{1}{2}n^2-15\,n^{3/2}.
\]
\end{theorem}

\begin{proposition}[Density of ones]\label{prop:density}
\leavevmode
\begin{enumerate}[label=\normalfont(\roman*)]
\item There exists an explicit family of strings $s_n$ with
  $f(s_n)\ge\tfrac{1}{2}n^2-15n^{3/2}$ and $\zeros(s_n)=O(\sqrt{n})$,
  so the ones-density satisfies $\ones(s_n)/n=1-O(n^{-1/2})$.
\item Every optimal string $s_n^*$ satisfies $\zeros(s_n^*)=O(n^{3/4})$,
  equivalently $\ones(s_n^*)/n=1-O(n^{-1/4})$.
\end{enumerate}
\end{proposition}

\begin{theorem}[Convexity]\label{thm:convex}
For every $n\ge2$,
\[
  a(n+1)-2a(n)+a(n-1)\;\ge\;0.
\]
Equivalently, $\Delta a(n):=a(n)-a(n-1)$ is nondecreasing for $n\ge2$.
\end{theorem}

Section~\ref{sec:upper} gives the upper bound.
Sections~\ref{sec:construction}--\ref{sec:lower} prove
Theorems~\ref{thm:limit}--\ref{thm:quant}.
Section~\ref{sec:density} proves Proposition~\ref{prop:density}.
Sections~\ref{sec:zeros}--\ref{sec:runs} develop the structural bounds.
Section~\ref{sec:convex} proves Theorem~\ref{thm:convex}.

\section{Definitions}

If $w=w_1\cdots w_k\in\{0,1\}^k$, its binary value is
$\val(w)=\sum_{i=1}^k w_i\,2^{k-i}$.
The value set of $s\in\{0,1\}^n$ is
$V(s)=\{\val(s_i\cdots s_j):1\le i\le j\le n\}$,
and $f(s)=|V(s)|$.
Write $\ones(s)$ and $\zeros(s)$ for the number of ones and zeros in $s$.

\section{Upper bound}\label{sec:upper}

\begin{lemma}[Upper bound]\label{lem:upper}
For every $n\ge1$, $a(n)\le n(n+1)/2$, and hence
$\limsup_{n\to\infty}a(n)/n^2\le1/2$.
\end{lemma}

\begin{proof}
Every distinct positive integer in $V(s)$ has a unique standard binary
representation beginning with~$\mathtt{1}$, so it is represented by at
least one substring of $s$ beginning with~$\mathtt{1}$.  Different
positive integers require different such substrings.

If $s$ contains at least one zero then $0\in V(s)$, represented only by
substrings beginning with~$\mathtt{0}$.  At least one of the $n(n+1)/2$
substrings begins with~$\mathtt{0}$, so the number of leading-$\mathtt{1}$
substrings is at most $n(n+1)/2-1$, giving $f(s)\le n(n+1)/2$.  If $s$
is all ones then $f(s)$ counts distinct positive values, at most $n(n+1)/2$.
\end{proof}

\begin{remark}
The bound is tight at $n=2$: $\mathtt{10}$ gives $V(\mathtt{10})=\{0,1,2\}$,
so $a(2)=3=2\cdot3/2$.
\end{remark}

\section{The construction}\label{sec:construction}

\subsection{Strictly decreasing partitions}

\begin{lemma}[Decreasing partition]\label{lem:partition}
Let $k\ge2$ and $M\ge k(k+1)/2$.  Then there exist strictly decreasing
positive integers $b_1>b_2>\cdots>b_k\ge1$ with $\sum_{i=1}^k b_i=M$.
Moreover, they may be chosen with $|b_i-M/k|\le(k+1)/2$ for all $i$.
\end{lemma}

\begin{proof}
\textit{Strategy.}  Consider the arithmetic sequence $c_i=q+k-i$
($i=1,\ldots,k$), which decreases by~$1$ at each step.  Its sum is
\begin{equation}\label{eq:basesum}
  \sum_{i=1}^{k}c_i
  = kq + \sum_{i=1}^{k}(k-i)
  = kq + \frac{k(k-1)}{2}.
\end{equation}
Choose $q$ so this base sum is as large as possible without exceeding
$M$, and let $r=M-\sum c_i$ be the shortfall.  We then add $1$ to the
first $r$ terms to account for the shortfall.

\textit{Definitions.}  Set
\[
  q = \Bigl\lfloor \frac{M}{k} - \frac{k-1}{2} \Bigr\rfloor,
  \qquad
  r = M - kq - \frac{k(k-1)}{2},
\]
and define
\[
  b_i=\begin{cases}q+k-i+1,&1\le i\le r,\\q+k-i,&r<i\le k.\end{cases}
\]

\textit{Step~1: $0\le r<k$.}
From the definition of $q$ as a floor, $q\le M/k-(k-1)/2<q+1$.
Multiplying through by $k$ gives
\[
  kq \;\le\; M - \tfrac{k(k-1)}{2} \;<\; k(q+1),
\]
so $0\le M-kq-k(k-1)/2<k$, i.e.\ $0\le r<k$.

\textit{Step~2: Strict decrease.}
We check all three cases for consecutive pairs $b_i>b_{i+1}$.
\begin{itemize}[nosep]
  \item \textit{Both in the first case} ($i<r$, so $i+1\le r$ as well):
    $b_i-b_{i+1}=(q+k-i+1)-(q+k-(i+1)+1)=1>0$.
  \item \textit{Transition} ($i=r\ge1$, so $i+1>r$):
    $b_r=q+k-r+1$ and $b_{r+1}=q+k-(r+1)$,
    giving $b_r-b_{r+1}=q+k-r+1-(q+k-r-1)=2>0$.
  \item \textit{Both in the second case} ($i>r$):
    $b_i-b_{i+1}=(q+k-i)-(q+k-i-1)=1>0$.
\end{itemize}
Hence $b_1>b_2>\cdots>b_k$.

\textit{Step~3: The sum equals $M$.}
Split the sum at the boundary between the two cases:
\begin{align*}
  \sum_{i=1}^{k}b_i
  &= \sum_{i=1}^{r}(q+k-i+1) + \sum_{i=r+1}^{k}(q+k-i) \\
  &= \sum_{i=1}^{k}(q+k-i) + \sum_{i=1}^{r}1 \\
  &= kq + \frac{k(k-1)}{2} + r \;=\; M,
\end{align*}
where the last equality uses the definition of~$r$ and~\eqref{eq:basesum}.

\textit{Step~4: All terms are positive.}
From $M\ge k(k+1)/2$ we get $M/k\ge(k+1)/2$, so
\[
  \frac{M}{k}-\frac{k-1}{2}\;\ge\;\frac{k+1}{2}-\frac{k-1}{2}=1,
\]
hence $q=\lfloor M/k-(k-1)/2\rfloor\ge1$.
Since $r<k$, the index $i=k$ falls in the second case, giving
$b_k=q+k-k=q\ge1$.

\textit{Step~5: Proximity bound.}
Every first-case term satisfies $q < q+k-r+1\le b_i\le q+k$, and
every second-case term satisfies $q\le b_i=q+k-i\le q+k-r-1<q+k$.
So all $b_i\in[q,\,q+k]$.
From the floor definition, $M/k\in[q+(k-1)/2,\;q+(k+1)/2)$,
and the width of $[q,q+k]$ around any point in this interval is at most
$(k+1)/2$ in each direction.
Hence $|b_i-M/k|\le(k+1)/2$ for all~$i$.
\end{proof}

\subsection{Parameters and string}

Fix $n\ge64$ and set
\begin{equation}\label{eq:params}
  k=\floor{\tfrac{\sqrt{n}}{4}},\qquad M=n-(k-1).
\end{equation}
Since $n\ge64$ we have $k=\lfloor\sqrt{n}/4\rfloor\ge\lfloor8/4\rfloor=2$.

\textit{Why $M\ge n/2$.}
Since $k\le\sqrt{n}/4$, we have $\sqrt{n}/4\le n/2$ (equivalently $\sqrt{n}\le2n$, true for all $n\ge1$),
so $k-1<k\le\sqrt{n}/4\le n/2$.
Hence $M=n-(k-1)>n-n/2=n/2$.

\textit{Why $n/2\ge k(k+1)/2$.}
This is equivalent to $n\ge k(k+1)$.  Since $k\le\sqrt{n}/4$ we have
\[
  k(k+1)\le k\cdot2k=2k^2\le2\cdot\frac{n}{16}=\frac{n}{8}\le n,
\]
where the second inequality uses $k\le\sqrt{n}/4$ so $k^2\le n/16$.
Hence $M\ge n/2\ge k(k+1)/2$, and Lemma~\ref{lem:partition} applies.

Choose $b_1>b_2>\cdots>b_k\ge1$ with $\sum b_i=M$ and
$|b_i-M/k|\le k$, as in Lemma~\ref{lem:partition}.  Define
\begin{equation}\label{eq:string}
  s\;=\;\underbrace{1^{b_1}}_{\text{block }1}\,0\,
         \underbrace{1^{b_2}}_{\text{block }2}\,0\,
         \cdots\,0\,
         \underbrace{1^{b_k}}_{\text{block }k},
\end{equation}
consisting of $k$ one-blocks separated by $k-1$ single zeros.  Its
length is $\sum b_i+(k-1)=M+(k-1)=n$.

\section{An injective family of substrings}\label{sec:family}

For integers $i,j,\alpha,\beta$ satisfying
\begin{equation}\label{eq:range}
  1\le i,\quad j\ge i+2,\quad j\le k,\quad
  1\le\alpha\le b_i,\quad 1\le\beta\le b_j,
\end{equation}
define $u_{i,j,\alpha,\beta}$ to be the substring of $s$ beginning
$\alpha$ positions before the end of block~$i$ and ending after the
first $\beta$ ones of block~$j$.  In binary,
\begin{equation}\label{eq:u}
  u_{i,j,\alpha,\beta}
  =1^{\alpha}\,0\,1^{b_{i+1}}\,0\,1^{b_{i+2}}\,0\,
  \cdots\,1^{b_{j-1}}\,0\,1^{\beta}.
\end{equation}
The condition $j\ge i+2$ ensures at least one full interior block
($b_{i+1},\ldots,b_{j-1}$).  Every $u_{i,j,\alpha,\beta}$ starts and
ends with~$\mathtt{1}$, so it is already a standard binary expansion of
a positive integer.

\begin{lemma}[Injectivity]\label{lem:inject}
The map $(i,j,\alpha,\beta)\mapsto u_{i,j,\alpha,\beta}$ is injective
on the range~\eqref{eq:range}.
\end{lemma}

\begin{proof}
Suppose $u_{i,j,\alpha,\beta}=u_{i',j',\alpha',\beta'}$ as binary words.
The run-length encoding (RLE) of~\eqref{eq:u}, listing alternating
run-lengths of ones then zeros, is
\[
  (\alpha,\,1,\,b_{i+1},\,1,\,b_{i+2},\,1,\,\ldots,\,b_{j-1},\,1,\,\beta).
\]
Matching entries gives $\alpha=\alpha'$, $\beta=\beta'$, and
$(b_{i+1},\ldots,b_{j-1})=(b_{i'+1},\ldots,b_{j'-1})$.  Since
$b_1>\cdots>b_k$ is strictly decreasing, every consecutive subsequence
occurs in exactly one position, so $i=i'$ and $j=j'$.
\end{proof}

\section{Quantitative lower bound and the limit}\label{sec:lower}

\begin{proof}[Proof of Theorem~\ref{thm:quant}]
By Lemma~\ref{lem:inject}, the family $\{u_{i,j,\alpha,\beta}\}$
consists of pairwise distinct positive integers, so
\begin{equation}\label{eq:f-lower}
  a(n)\ge f(s)\ge\sum_{\substack{1\le i<j\le k\\j\ge i+2}}b_i b_j.
\end{equation}
Expanding:
\begin{equation}\label{eq:expand}
  \sum_{j\ge i+2}b_i b_j
  =\frac{(\sum b_i)^2-\sum b_i^2}{2}-\sum_{i=1}^{k-1}b_i b_{i+1}.
\end{equation}

\textbf{Step~1.}  Since $\sum b_i=M=n-k+1$ and $k\le\sqrt{n}/4$,
\[
  \frac{M^2}{2}\ge\frac{(n-\sqrt{n}/4)^2}{2}
  \ge\frac{n^2}{2}-\frac{n^{3/2}}{4}.
\]

\textbf{Step~2.}  Write $b_i=M/k+d_i$ with $|d_i|\le k$.  Then
$\sum b_i^2\le M^2/k+k^3$.
Since $k\ge\sqrt{n}/8$ (as $\lfloor x\rfloor\ge x/2$ for $x\ge2$),
$M^2/k\le n^2/(\sqrt{n}/8)=8n^{3/2}$,
and $k^3\le(\sqrt{n}/4)^3=n^{3/2}/64$, so
\begin{equation}\label{eq:sumsq}
  \sum_{i=1}^k b_i^2<9n^{3/2}.
\end{equation}

\textbf{Step~3.}  By AM--GM,
$b_i b_{i+1}\le(b_i^2+b_{i+1}^2)/2$, so
\begin{equation}\label{eq:adj}
  \sum_{i=1}^{k-1}b_i b_{i+1}\le\sum_{i=1}^k b_i^2<9n^{3/2}.
\end{equation}

\textbf{Conclusion.}  Substituting into~\eqref{eq:expand}:
\[
  a(n)\ge\frac{n^2}{2}-\frac{n^{3/2}}{4}
  -\frac{9n^{3/2}}{2}-9n^{3/2}
  =\frac{n^2}{2}-\frac{55}{4}n^{3/2}>\frac{n^2}{2}-15n^{3/2}.\qedhere
\]
\end{proof}

\begin{proof}[Proof of Theorem~\ref{thm:limit}]
Lemma~\ref{lem:upper} gives $\limsup a(n)/n^2\le1/2$.
Theorem~\ref{thm:quant} gives $a(n)/n^2\ge1/2-15/\sqrt{n}\to1/2$.
\end{proof}

\begin{corollary}\label{cor:rate}
For all $n\ge64$, $\;n^2/2-15n^{3/2}\le a(n)\le n(n+1)/2$, so
$a(n)/n^2=1/2+O(n^{-1/2})$.
\end{corollary}

\begin{remark}
The lower bound $n^2/2-15n^{3/2}$ is negative for $n<900$ (since
$n^2/2>15n^{3/2}$ iff $\sqrt{n}>30$).  For such $n$ the inequality in
Theorem~\ref{thm:quant} is trivially satisfied by the non-negativity of
$a(n)$, and carries no information.  The bound is meaningful only for
$n\ge900$; for small~$n$ the exact values from computational search
(Table~\ref{tab:evidence}) give much sharper information.
\end{remark}

\subsection{Sharper quantitative lower bound}\label{sec:quant-sharp}

The constant $15$ in Theorem~\ref{thm:quant} is an artefact of the
parameter choice $k=\lfloor\sqrt n/4\rfloor$.  Choosing instead
$k=\lceil\sqrt n\,\rceil$, together with a slightly sharper bound on
the variance of the partition produced by Lemma~\ref{lem:partition},
reduces the constant by a factor of $60/13>4.6$, and makes the bound
non-trivial already at $n=36$ rather than $n=900$.

\begin{theorem}[Sharper quantitative lower bound]\label{thm:quant-sharp}
For every $n\ge 36$,
\[
  a(n)\;\ge\;\frac{n^2}{2}-\frac{13}{4}\,n^{3/2}.
\]
\end{theorem}

\begin{remark}\label{rem:quant-sharp-comparison}
Comparison with Theorem~\ref{thm:quant}:
\begin{itemize}[nosep]
  \item The new bound is positive iff $\sqrt n>13/2$, i.e.\ $n\ge 43$,
    while the old bound is positive only for $n\ge 901$.  Thus the
    threshold beyond which the bound is informative drops by a factor
    of more than~$20$.
  \item The constant in front of $n^{3/2}$ improves from $15$ to
    $13/4$, a factor of $60/13>4.6$.  In fact the construction
    achieves a constant of $\sqrt 6+o(1)\approx 2.449$
    (Remark~\ref{rem:quant-optimal} below); we keep $13/4$ here
    because it admits a fully symbolic and elementary proof.
\end{itemize}
A numerical comparison at sample $n$ (with $a(n)$ taken from
the OEIS b-file b112509):
\begin{center}\small
  \begin{tabular}{r r r r}
    \toprule
    $n$    & Thm.~\ref{thm:quant} & Thm.~\ref{thm:quant-sharp} & $a(n)$ \\
    \midrule
    $36$   & $-2790$   & $-54$   & $450$ \\
    $50$   & $-3050$   & $101$   & $899$ \\
    $100$  & $-10000$  & $1750$  & $3875$ \\
    $150$  & $-16320$  & $5278$  & $9070$ \\
    $1000$ & $-237201$ & $397214$ & --- \\
    \bottomrule
  \end{tabular}
\end{center}
\end{remark}

\begin{proof}[Proof of Theorem~\ref{thm:quant-sharp}]
The construction~\eqref{eq:string} and Lemmas~\ref{lem:partition} and
\ref{lem:inject} are unchanged.  Only the parameter choice and the
estimation of the loss terms in~\eqref{eq:expand} differ.

\smallskip\noindent\textit{Parameter.}  Set
$k=\lceil\sqrt n\,\rceil$ and $M=n-k+1$.  Then
$\sqrt n\le k\le\sqrt n+1$.

\smallskip\noindent\textit{Lemma~\ref{lem:partition} applies for
$n\ge 36$.}  We need $k\ge 2$ (clear) and $M\ge k(k+1)/2$.  The
latter is equivalent to $n-\sqrt n\ge\tfrac12(\sqrt n+1)(\sqrt n+2)
=\tfrac{n}{2}+\tfrac{3\sqrt n}{2}+1$ (using $k\le\sqrt n+1$), i.e.\
$\tfrac{n}{2}\ge\tfrac{5\sqrt n}{2}+1$, which holds for
$\sqrt n\ge 6$, i.e.\ $n\ge 36$.

\smallskip\noindent\textit{Variance bound.}
By Lemma~\ref{lem:partition}, every $b_i$ lies in $[q,q+k]$ where
$q=\lfloor M/k\rfloor$.  Since $\sum_{i=1}^k(b_i-M/k)=0$ and each
deviation has absolute value at most $k/2+1\le(k+1)/2$
(because $b_i\in\{q,\ldots,q+k\}$ and the average is $M/k\in[q,q+k]$),
\begin{equation}\label{eq:sumsq-sharp}
  \sum_{i=1}^k b_i^2
  \;=\;\frac{M^2}{k}+\sum_{i=1}^k\!\Bigl(b_i-\tfrac{M}{k}\Bigr)^{\!2}
  \;\le\;\frac{M^2}{k}+k\cdot\frac{(k+1)^2}{4}.
\end{equation}
By the AM--GM inequality
$2b_ib_{i+1}\le b_i^2+b_{i+1}^2$,
\begin{equation}\label{eq:adjprod-sharp}
  \sum_{i=1}^{k-1}b_ib_{i+1}
  \;\le\;\sum_{i=1}^k b_i^2
  \;\le\;\frac{M^2}{k}+\frac{k(k+1)^2}{4}.
\end{equation}
Substituting~\eqref{eq:sumsq-sharp} and~\eqref{eq:adjprod-sharp}
into~\eqref{eq:expand}:
\begin{align*}
  a(n)
  &\;\ge\;\frac{M^2-\sum b_i^2}{2}-\sum_{i=1}^{k-1}b_ib_{i+1} \\
  &\;\ge\;\frac{M^2}{2}-\frac{1}{2}\!\left(\frac{M^2}{k}+\frac{k(k+1)^2}{4}\right)
        -\!\left(\frac{M^2}{k}+\frac{k(k+1)^2}{4}\right) \\
  &\;=\;\frac{M^2}{2}-\frac{3M^2}{2k}-\frac{3k(k+1)^2}{8}.
\end{align*}

\smallskip\noindent\textit{Term-by-term estimates.}
Using $\sqrt n\le k\le\sqrt n+1$ and $M=n-k+1$:
\begin{enumerate}[label=(\roman*),nosep]
\item $\dfrac{M^2}{2}\ge\dfrac{(n-\sqrt n)^2}{2}
  =\dfrac{n^2}{2}-n^{3/2}+\dfrac{n}{2}$.
\item $\dfrac{3M^2}{2k}\le\dfrac{3n^2}{2\sqrt n}=\dfrac{3}{2}n^{3/2}$.
\item Expanding $(\sqrt n+1)(\sqrt n+2)^2=n^{3/2}+5n+8\sqrt n+4$
  gives
  \[
    \frac{3k(k+1)^2}{8}
    \;\le\;\frac{3(\sqrt n+1)(\sqrt n+2)^2}{8}
    \;=\;\frac{3}{8}n^{3/2}+\frac{15}{8}n+3\sqrt n+\frac{3}{2}.
  \]
\end{enumerate}
Combining (i)--(iii):
\begin{align}
  a(n)
  &\;\ge\;\frac{n^2}{2}-n^{3/2}+\frac{n}{2}
        -\frac{3}{2}n^{3/2}
        -\frac{3}{8}n^{3/2}-\frac{15}{8}n-3\sqrt n-\frac{3}{2}
        \nonumber\\
  &\;=\;\frac{n^2}{2}-\frac{23}{8}n^{3/2}-\frac{11}{8}n-3\sqrt n-\frac{3}{2}.
        \label{eq:lb-allterms}
\end{align}

\smallskip\noindent\textit{Tail bound.}  We claim
\begin{equation}\label{eq:tail}
  \frac{11}{8}n+3\sqrt n+\frac{3}{2}\;\le\;\frac{3}{8}n^{3/2}
  \qquad(n\ge 36).
\end{equation}
Equivalently $11n+24\sqrt n+12\le 3n^{3/2}$.  At $n=36$ the right
side equals $3\cdot 216=648$ and the left side equals
$396+144+12=552$, so the inequality holds at $n=36$ with margin
$96$.  Differentiating, $\tfrac{d}{dn}(3n^{3/2}-11n-24\sqrt n-12)
=\tfrac{9}{2}\sqrt n-11-\tfrac{12}{\sqrt n}$, which is positive for
$n\ge 36$ since $\tfrac{9}{2}\cdot 6-11-2=14>0$ and the
derivative is monotone increasing.  Therefore the inequality holds
for all $n\ge 36$.

Combining~\eqref{eq:lb-allterms} and~\eqref{eq:tail},
\[
  a(n)\;\ge\;\frac{n^2}{2}-\frac{23}{8}n^{3/2}-\frac{3}{8}n^{3/2}
       \;=\;\frac{n^2}{2}-\frac{26}{8}n^{3/2}
       \;=\;\frac{n^2}{2}-\frac{13}{4}n^{3/2}. \qedhere
\]
\end{proof}

\begin{remark}\label{rem:quant-optimal}
The constant $13/4$ is not optimal.  Tracking the variance of the
partition more carefully (it is bounded by $k(k^2-1)/12+O(k)$ rather
than $k(k+1)^2/4$) and choosing $k=\lceil\sqrt{3n/2}\,\rceil$ instead
of $\lceil\sqrt n\,\rceil$ yields
\[
  a(n)\;\ge\;\frac{n^2}{2}-\sqrt 6\,n^{3/2}+O(n),
\]
with $\sqrt 6\approx 2.449$.  The choice $k=\lceil\sqrt{3n/2}\,\rceil$
is the one that equates the two leading $n^{3/2}$ losses $kn$ and
$3n^2/(2k)$ in the proof above.  We do not record the technical
details since they would obscure the elementary nature of
Theorem~\ref{thm:quant-sharp}.
\end{remark}

\begin{corollary}\label{cor:rate-sharp}
For all $n\ge 36$, $\;a(n)/n^2\;=\;\tfrac12+O(n^{-1/2})$ with explicit
constants:
\(
  \tfrac12-\tfrac{13}{4}n^{-1/2}\;\le\; a(n)/n^2\;\le\;\tfrac12+\tfrac{1}{2n}.
\)
\end{corollary}

\section{A structural constraint on optima}\label{sec:structure}

In this section we prove a rigorous upper bound on $f(s)$ in terms
of the run-length statistics of $s$.  The bound is universal: it
holds for \emph{every} binary string and is tight enough to force
quantitative structural constraints on every optimum of $a(n)$.

\subsection{Run-length notation}

For $s\in\{0,1\}^n$, write its \emph{run decomposition}
\[
  s\;=\;a^{r_1}\,\bar a^{r_2}\,a^{r_3}\,\cdots
\]
where $a\in\{0,1\}$ is the leading symbol and the maximal runs have
lengths $r_1,r_2,\ldots$ alternating between $a$ and~$\bar a$.  Let
$\ell_1,\ldots,\ell_p$ denote the lengths of the $1$-runs and
$m_1,\ldots,m_q$ the lengths of the $0$-runs (so
$p+q\le n+1$ and $\sum\ell_t+\sum m_t=n$).  Set
\[
  L(s)=\max_t\ell_t,\qquad K(s)=\max_t m_t,
\]
the \emph{longest $1$-run} and \emph{longest $0$-run} respectively
(with the convention $L(s)=0$ if $s$ contains no $1$s, and
$K(s)=0$ if $s$ contains no $0$s).

\subsection{The run-length frequency bound}

We work with two versions of the bound: a tighter \emph{asymmetric}
form involving only the longest $1$-run, and a slightly looser
\emph{symmetric} form that lets us simultaneously control both the
longest $1$-run and the longest $0$-run.

\begin{theorem}[Run-length frequency bound, asymmetric]\label{thm:rlf}
Every $s\in\{0,1\}^n$ satisfies
\begin{equation}\label{eq:rlf}
  f(s)\;\le\;1+L(s)+\binom{n+1}{2}
        -\sum_{t=1}^{p}\binom{\ell_t+1}{2}
        -\sum_{t=1}^{q}\binom{m_t+1}{2}.
\end{equation}
\end{theorem}

\begin{proof}
Recall $f(s)=|V(s)|$, where $V(s)\subseteq\mathbb Z_{\ge 0}$ is the
set of integer values represented by substrings of $s$, with the
empty substring contributing the value $0$.  Partition
\begin{equation}\label{eq:partition}
  V(s)\;=\;\{0\}\,\cup\,V_1(s)\,\cup\,V_{\mathrm m}(s),
\end{equation}
where $V_1(s)$ is the set of values represented by some all-$1$
substring of $s$ (necessarily $\ne 0$), and
$V_{\mathrm m}(s):=V(s)\setminus(\{0\}\cup V_1(s))$ collects the
remaining values.  This is a disjoint union: $0\notin V_1$ since
every length-$\ge 1$ all-$1$ substring has value $2^j-1\ge 1$, and
$V_{\mathrm m}$ is defined to be disjoint from the other two sets.

\smallskip\noindent\textit{Counting $V_1(s)$.}
A nonempty substring of the form $0^a 1^j$ ($a\ge 0$, $j\ge 1$) has
integer value $2^j-1$, depending only on $j$.  Hence
$V_1(s)=\{2^j-1:1\le j\le L(s)\}$ and $|V_1(s)|=L(s)$.

\smallskip\noindent\textit{Counting $V_{\mathrm m}(s)$.}
Every value in $V_{\mathrm m}(s)$ is represented by at least one
substring, and that substring must contain both a $0$ and a $1$:
otherwise it would be a power of $2$ minus one (hence in $V_1$) or
zero.  Therefore $|V_{\mathrm m}(s)|$ is bounded above by the number
of pairs $(i,j)$ with $1\le i\le j\le n$ such that
$s_i\cdots s_j$ contains both symbols.  This count equals
$\binom{n+1}{2}$ minus the number of $(i,j)$ for which
$s_i\cdots s_j$ is all-$1$ or all-$0$.  A pair $(i,j)$ yields an
all-$1$ substring iff $\{i,\ldots,j\}$ is contained in some maximal
$1$-run, and the number of such pairs within a run of length
$\ell$ is $\binom{\ell+1}{2}$.  Summing over $1$-runs and $0$-runs,
\[
  |V_{\mathrm m}(s)|\;\le\;
  \binom{n+1}{2}-\sum_{t=1}^{p}\binom{\ell_t+1}{2}
                -\sum_{t=1}^{q}\binom{m_t+1}{2}.
\]
Combining with $|V_1(s)|=L(s)$ via~\eqref{eq:partition}
gives~\eqref{eq:rlf}.
\end{proof}

\begin{remark}\label{rem:rlf-tight}
Inequality~\eqref{eq:rlf} has been verified to hold with strictly
positive slack for every optimal string in our database
($1\le n\le 150$, totalling $21\,072$ optima).  At sample~$n$ the
slack $\mathrm{RHS}-a(n)$ behaves like:
\begin{center}\small
\begin{tabular}{r r r r r r r r}
\toprule
$n$ & 10 & 30 & 50 & 80 & 100 & 130 & 150 \\
\midrule
slack & 7 & 77 & 176 & 436 & 614 & 965 & 1\,241 \\
\bottomrule
\end{tabular}
\end{center}
matching the predicted $\Theta(n^{3/2})$ scaling
of the gap $\binom{n+1}{2}-a(n)$.  The bound becomes equality when
every pair $(i,j)$ yielding a mixed substring gives an integer value
not already counted; this fails for long $s$ because of the
leading-zero collision $0^a 1^j 0^b\sim 1^j 0^b$.
\end{remark}

\subsection{Forced run-length statistics at any optimum}

\begin{corollary}[Concentration of run lengths]\label{cor:concentration}
Let $s$ be any optimum for $a(n)$ with run lengths
$\ell_1,\ldots,\ell_p,m_1,\ldots,m_q$ as above and longest
$1$-run $L=L(s)$.  Then for all $n\ge 36$,
\begin{equation}\label{eq:concentration}
  \sum_{t=1}^{p}\binom{\ell_t+1}{2}
  +\sum_{t=1}^{q}\binom{m_t+1}{2}
  \;\le\;\frac{13}{4}\,n^{3/2}+\frac{n}{2}+L+1.
\end{equation}
\end{corollary}

\begin{proof}
Apply Theorem~\ref{thm:rlf} with $f(s)=a(n)$ and rearrange:
\[
  \sum_t\binom{\ell_t+1}{2}+\sum_t\binom{m_t+1}{2}
  \;\le\;\binom{n+1}{2}+1+L-a(n).
\]
Using $\binom{n+1}{2}=\tfrac{n^2}{2}+\tfrac{n}{2}$ and the lower
bound $a(n)\ge\tfrac{n^2}{2}-\tfrac{13}{4}n^{3/2}$ of
Theorem~\ref{thm:quant-sharp},
$\binom{n+1}{2}-a(n)\le\tfrac{n}{2}+\tfrac{13}{4}n^{3/2}$, giving
\eqref{eq:concentration}.
\end{proof}

\begin{corollary}[Forced upper bound on the longest $1$-run]\label{cor:Lbound}
For every $n\ge 36$ and every optimum $s$ for $a(n)$,
\begin{equation}\label{eq:Lbound}
  \binom{L(s)+1}{2}
  \;\le\;\frac{13}{4}\,n^{3/2}+\frac{n}{2}+L(s)+1,
\end{equation}
hence
$L(s)\;\le\;\sqrt{\tfrac{13}{2}}\;n^{3/4}+O(n^{3/8})\;=\;O(n^{3/4})$.
\end{corollary}

\begin{proof}
The longest $1$-run contributes $\binom{L+1}{2}$ to the left-hand
side of~\eqref{eq:concentration}; all other terms are non-negative,
giving~\eqref{eq:Lbound}.  Since $\binom{L+1}{2}\ge L^2/2$, the
inequality forces $L^2/2\le\tfrac{13}{4}n^{3/2}+L+\tfrac{n}{2}+1$.
Solving and absorbing lower-order terms,
$L\le\sqrt{13/2}\,n^{3/4}+O(n^{3/8})$.
\end{proof}

\subsection{A symmetric bound and the longest $0$-run}

The argument above privileges $1$-substrings because the empty
substring shares the value $0$ with every all-$0$ substring,
making the all-$0$ substrings ``invisible'' to $f$.  By treating
$0$-substrings on equal footing one obtains a slightly looser
bound that controls both the longest $1$-run and the longest
$0$-run.

\begin{theorem}[Run-length frequency bound, symmetric]\label{thm:rlf-sym}
Every $s\in\{0,1\}^n$ satisfies
\begin{equation}\label{eq:rlf-sym}
  f(s)\;\le\;1+L(s)+K(s)+\binom{n+1}{2}
        -\sum_{t=1}^{p}\binom{\ell_t+1}{2}
        -\sum_{t=1}^{q}\binom{m_t+1}{2}.
\end{equation}
\end{theorem}

\begin{proof}
Inequality~\eqref{eq:rlf-sym} follows immediately from
Theorem~\ref{thm:rlf} since $K(s)\ge 0$.
\end{proof}

\begin{corollary}[Forced upper bound on the longest $0$-run]\label{cor:Kbound}
For every $n\ge 36$ and every optimum $s$ for $a(n)$,
\[
  \binom{K(s)+1}{2}
  \;\le\;\frac{13}{4}\,n^{3/2}+\frac{n}{2}+L(s)+K(s)+1,
\]
hence $K(s)=O(n^{3/4})$.
\end{corollary}

\begin{proof}
Apply Theorem~\ref{thm:rlf-sym} and rearrange exactly as in the
proof of Corollary~\ref{cor:concentration}, retaining the
$\binom{K+1}{2}$ term on the left-hand side and dropping the
non-negative remaining run terms.
\end{proof}

\subsection{Empirical refinement}

It is tempting to conjecture that the leading $1$-run length $b_1(n)$
of an optimum is monotone non-decreasing in $n$.  This is
\emph{false}: in the database of all optima for $1\le n\le 150$ we
find exactly three values of $n$ at which the minimum leading
$1$-run length \emph{decreases}, namely
\[
  b_1(51)=12<b_1(50)=13,\quad
  b_1(71)=15<b_1(70)=16,\quad
  b_1(116)=20<b_1(115)=21.
\]
However the slightly weaker monotonicity
\[
  \min_{s\in S_n}b_1(s)\;\ge\;\min_{s\in S_{n-4}}b_1(s)
\]
holds for every $n\le 150$ in our database, where $S_n$ is the set of
optima for length~$n$.  We do not have a proof of this gap-$4$
monotonicity; the run-length frequency
bound~\eqref{eq:rlf} only constrains $L(s)=O(n^{3/4})$ from above
and is silent on monotonicity.

The construction of \S\ref{sec:construction} achieves leading
$1$-run length $b_1=q+k=\Theta(\sqrt n)$, well below the
information-theoretic upper bound $O(n^{3/4})$ of
Corollary~\ref{cor:Lbound}.  Whether every optimum satisfies
$b_1(s)=\Theta(\sqrt n)$ (as the data suggest) is an open problem.

\section{Density of ones}\label{sec:density}

\subsection{Near-optimal strings}

\begin{proof}[Proof of Proposition~\ref{prop:density}(i)]
The string $s=s_n$ from~\eqref{eq:string} satisfies
$f(s_n)\ge n^2/2-15n^{3/2}$ (this is established directly in the proof
of Theorem~\ref{thm:quant}, as the lower bound on $a(n)$ is obtained by
exhibiting $s_n$ and bounding $f(s_n)$ from below).

\textit{Counting the zeros.}
By construction, $s_n$ consists of $k$ one-blocks separated by $k-1$
single zeros, so it contains exactly $k-1$ zeros.  Since
$k=\lfloor\sqrt{n}/4\rfloor\le\sqrt{n}/4$, we have
\[
  \zeros(s_n) = k-1 \le \frac{\sqrt{n}}{4}-1 < \frac{\sqrt{n}}{4},
\]
so $\zeros(s_n)=O(\sqrt{n})$.

\textit{Deducing the ones-density.}
Since $s_n$ has length $n$,
\[
  \ones(s_n) = n - \zeros(s_n) \ge n - \frac{\sqrt{n}}{4},
\]
and therefore
\[
  \frac{\ones(s_n)}{n} \ge 1 - \frac{1}{4\sqrt{n}} = 1 - O(n^{-1/2}).\qedhere
\]
\end{proof}

\subsection{Zero-position penalty}

\begin{lemma}[Zero-position bound]\label{lem:zero-pos}
Let $s\in\{0,1\}^n$ have zeros at positions $p_1<\cdots<p_z$.  Then
\[
  f(s)\le\frac{n(n+1)}{2}-\sum_{r=1}^z(n-p_r+1).
\]
In particular, $\;f(s)\le n(n+1)/2-z(z+1)/2$.
\end{lemma}

\begin{proof}
For each $x\in V(s)$ choose a minimum-length representative $w_x$.
The map $x\mapsto w_x$ is injective.  If $x>0$ then $w_x$ starts with
$\mathtt{1}$ (otherwise removing leading zeros gives a shorter
representative).  So distinct positive values correspond to distinct
leading-$\mathtt{1}$ substrings.

\textit{First inequality.}
A substring starting at position $p_r$ (which holds a zero) has the
form $0\,s_{p_r+1}\cdots s_j$ for some $j\ge p_r$, so it starts
with~$\mathtt{0}$.  The number of substrings that start at position
$p_r$ is $n-p_r+1$ (one for each ending position $j\in\{p_r,\ldots,n\}$).
None of these substrings starts with~$\mathtt{1}$, so none can be the
minimum-length representative $w_x$ of any positive integer $x$.
Removing all such substrings from consideration, the number of
available starting positions for positive-value representatives is at
most $n-z$ (the positions that hold a~$\mathtt{1}$), and the total
count of leading-$\mathtt{1}$ substrings is at most
$n(n+1)/2-\sum_{r=1}^{z}(n-p_r+1)$, giving the first inequality.

\textit{Second inequality.}
We minimise $\sum_{r=1}^{z}(n-p_r+1)$ subject to
$1\le p_1<p_2<\cdots<p_z\le n$.  Since each term $n-p_r+1$
is decreasing in $p_r$, the sum is minimised by taking the $p_r$ as
large as possible, i.e.\ $p_r=n-z+r$ for $r=1,\ldots,z$.
Substituting:
\[
  \sum_{r=1}^{z}(n-p_r+1)
  = \sum_{r=1}^{z}(n-(n-z+r)+1)
  = \sum_{r=1}^{z}(z-r+1)
  = \sum_{t=1}^{z}t
  = \frac{z(z+1)}{2},
\]
where we set $t=z-r+1$.  Hence the minimum of $\sum(n-p_r+1)$ is
$z(z+1)/2$, so $f(s)\le n(n+1)/2-z(z+1)/2$.
\end{proof}

\subsection{Optimal strings have few zeros}

\begin{proof}[Proof of Proposition~\ref{prop:density}(ii)]
Let $s_n^*$ be optimal with $z_n=\zeros(s_n^*)$.  By
Theorem~\ref{thm:quant} and Lemma~\ref{lem:zero-pos},
\[
  \frac{n^2}{2}-15n^{3/2}
  \le a(n)=f(s_n^*)
  \le\frac{n(n+1)}{2}-\frac{z_n(z_n+1)}{2}.
\]
Hence
\[
  \frac{z_n(z_n+1)}{2}
  \le\frac{n(n+1)}{2}-\frac{n^2}{2}+15n^{3/2}
  =\frac{n}{2}+15n^{3/2}.
\]
Multiplying through by $2$ and using $z_n(z_n+1)\ge z_n^2$:
\[
  z_n^2 \le z_n(z_n+1) \le n+30n^{3/2} \le 31n^{3/2}
  \quad\text{for all }n\ge1.
\]
Taking square roots, $z_n \le \sqrt{31}\,n^{3/4}$, so
$z_n=O(n^{3/4})$.  Since $\ones(s_n^*)=n-z_n$,
\[
  \frac{\ones(s_n^*)}{n} = 1-\frac{z_n}{n} \ge 1-\frac{\sqrt{31}}{n^{1/4}}
  = 1-O(n^{-1/4}).\qedhere
\]
\end{proof}

\section{Upper bounds via the zero structure}\label{sec:zeros}

Lemma~\ref{lem:zero-pos} treats all zeros equally.  We now refine the
bound using the \emph{run structure} of the zeros.

Let the maximal zero-runs of $s$ have lengths $c_1,\ldots,c_m$, with
the $i$-th run beginning at position $a_i$ and $r_i=n-(a_i+c_i)+1$
positions to its right.

\begin{lemma}[Run-sensitive bound]\label{lem:runs-raw}
$f(s)\le n(n+1)/2-\sum_{i=1}^m\bigl(c_i r_i+\binom{c_i+1}{2}\bigr)$.
\end{lemma}

\begin{proof}
Apply Lemma~\ref{lem:zero-pos} and group by zero-runs.
The $i$-th run occupies positions $a_i,\ldots,a_i+c_i-1$, contributing
$\sum_{t=0}^{c_i-1}(n-a_i-t+1)=c_i(n-a_i+1)-c_i(c_i-1)/2
=c_i r_i+\binom{c_i+1}{2}$ to the sum.
\end{proof}

\section{Fragmentation penalty}\label{sec:runs}

\begin{corollary}[Fragmentation]\label{cor:fragment}
If $s$ has $z$ zeros in $m$ maximal zero-runs, then
\[
  f(s)\;\le\;\frac{n(n+1)}{2}-\frac{z(z+1)}{2}-\binom{m}{2}.
\]
\end{corollary}

\begin{proof}
It suffices to show
$\sum_i\bigl(c_i r_i+\binom{c_i+1}{2}\bigr)\ge z(z+1)/2+\binom{m}{2}$.
Since consecutive zero-runs are separated by at least one $\mathtt{1}$,
for $i<m$ we have $r_i\ge\sum_{j>i}c_j+(m-i)$, so
$\sum_i c_i r_i\ge\sum_{i<j}c_i c_j+\sum_{i=1}^{m-1}c_i(m-i)$.
The identity $\binom{z+1}{2}=\sum_{i<j}c_i c_j+\sum_i(c_i^2+c_i)/2$ gives
$\sum_{i<j}c_i c_j+\sum_i\binom{c_i+1}{2}=z(z+1)/2$.
Finally $\sum_{i=1}^{m-1}c_i(m-i)\ge\sum_{i=1}^{m-1}(m-i)=\binom{m}{2}$
(each $c_i\ge1$).
\end{proof}

\begin{remark}
Corollary~\ref{cor:fragment} shows that fragmenting the zeros into more
runs always reduces the upper bound on $f(s)$, regardless of positions.
Each additional run incurs an extra penalty of at least $m-1$ (the new
run number), on top of the baseline $z(z+1)/2$.
\end{remark}

\section{Convexity of the sequence}\label{sec:convex}

The computational data (Table~\ref{tab:evidence} and~\cite{github})
show that the second differences
$\Delta^2 a(n):=a(n+1)-2a(n)+a(n-1)$ take only the values
$\{0,1,2\}$ for all $n\le200$.  We now prove one half of this
observation: $\Delta^2 a(n)\ge0$ for all $n\ge2$.  The proof is purely
combinatorial and does not rely on the explicit construction of
Section~\ref{sec:construction}.

\subsection{Background and strategy}

Write $\Delta a(n)=a(n)-a(n-1)$ for $n\ge2$.  The claim
$\Delta^2 a(n)\ge0$ is equivalent to
\begin{equation}\label{eq:convex-equiv}
  a(n+1)\;\ge\;a(n)+\Delta a(n)=2a(n)-a(n-1).
\end{equation}
That is, we must exhibit a string of length $n+1$ achieving at least
$a(n)+\Delta a(n)$ distinct substring values.

The strategy is to start from an \emph{optimal} string~$s$ of
length~$n$ and extend it by one bit (appending or prepending).  The key
insight is that the values in $V(s)$ that are ``new'' relative to a
length-$(n{-}1)$ truncation of~$s$ can each be ``lifted'' to a
genuinely new value in the extended string, and there are at least
$\Delta a(n)$ such values.

\subsection{Notation and setup}

Fix $n\ge2$ and let $s=s_1 s_2\cdots s_n$ be an optimal string of
length~$n$, so $f(s)=a(n)$.  Define the prefix and suffix of
length~$n-1$:
\begin{equation}\label{eq:pq}
  p=s_1 s_2\cdots s_{n-1},\qquad q=s_2 s_3\cdots s_n.
\end{equation}
Since every substring of~$p$ (resp.\ $q$) is also a substring of~$s$,
we have $V(p)\subseteq V(s)$ and $V(q)\subseteq V(s)$.

Since $p$ has length $n-1$, we have $f(p)\le a(n-1)$.  Therefore
\begin{equation}\label{eq:Vdiff-p}
  |V(s)\setminus V(p)| = f(s)-f(p)
    \;\ge\; a(n)-a(n-1)=\Delta a(n),
\end{equation}
and by the same argument applied to the suffix,
\begin{equation}\label{eq:Vdiff-q}
  |V(s)\setminus V(q)| \;\ge\; \Delta a(n).
\end{equation}
Each value $x\in V(s)\setminus V(p)$ is represented in~$s$ but not
in the prefix~$p$; every witness for~$x$ must ``touch'' the last
position.  Symmetrically, each value in $V(s)\setminus V(q)$ has every
witness touching the first position.

\subsection{Unique standard representation}

Recall that every positive integer $x$ has a \emph{standard binary
representation}: the unique binary string starting with~$\mathtt{1}$
and having value~$x$.  Its length is $\lfloor\log_2 x\rfloor+1$.
If $w$ is any other binary string with $\val(w)=x>0$, then either $w$
starts with~$\mathtt{0}$ (leading zeros) or has the same length as the
standard form; in the latter case it is the standard form itself.

For a positive $x\in V(s)$, define the \emph{canonical witness}
$w_x$ to be a minimum-length substring of~$s$ with $\val(w_x)=x$.
Then:

\begin{lemma}\label{lem:canon}
If $x>0$ then $w_x$ is the standard binary representation of~$x$.
\end{lemma}

\begin{proof}
If $w_x$ started with~$\mathtt{0}$, removing the leading zero gives a
strictly shorter substring of~$s$ with the same value, contradicting
the minimality of~$w_x$.  So $w_x$ starts with~$\mathtt{1}$.  Its
length is at least $\lfloor\log_2 x\rfloor+1$ (the length of the
standard form); by minimality, equality holds, so $w_x$ has the same
length as the standard form.  Since both start with~$\mathtt{1}$ and
represent~$x$, they are the same binary string.
\end{proof}

\begin{proof}[Proof of Theorem~\ref{thm:convex}]

\textbf{Step~1:\ new positive values must touch the boundary.}

Let $x\in V(s)\setminus V(p)$ with $x>0$.  Its canonical witness
$w_x$ is, by Lemma~\ref{lem:canon}, the standard representation
of~$x$; in particular $w_x$ starts with~$\mathtt{1}$.

\textit{Claim:}\ $w_x$ ends at position~$n$ (the last position
of~$s$).

\textit{Proof.}\ Write $w_x=s_i\cdots s_j$.  If $j\le n-1$ then
$w_x$ is a substring of $p=s_1\cdots s_{n-1}$, giving $x\in V(p)$
and contradicting $x\in V(s)\setminus V(p)$.  Hence $j=n$.

\medskip
\textbf{Step~2:\ appending a bit creates a new value.}

Fix any $b\in\{0,1\}$ and form the extended string
$sb=s_1\cdots s_n b$ of length $n+1$.

The word $w_x b$ occurs as a substring of~$sb$ ending at the new
last position.  Since $w_x$ starts with~$\mathtt{1}$ and has length
$\ell:=|w_x|$, the string $w_x b$ starts with~$\mathtt{1}$ and has
length $\ell+1$, representing the integer
$\val(w_x b)=2\,\val(w_x)+b=2x+b$.

\textit{Claim:}\ $w_x b$ is the standard representation of $2x+b$.

\textit{Proof.}\ Since $w_x$ is the standard representation of~$x$,
we have $2^{\ell-1}\le x<2^{\ell}$, so
$2^{\ell}\le 2x+b<2^{\ell+1}$.  Hence the standard representation
of~$2x+b$ has exactly $\ell+1$ bits.  The string $w_x b$ starts
with~$\mathtt{1}$, has $\ell+1$ bits, and has value $2x+b$; it is
therefore the standard representation.

\textit{Claim:}\ $2x+b\notin V(s)$.

\textit{Proof.}\ Suppose $2x+b\in V(s)$.  Take a minimum-length
representative~$u$ of~$2x+b$ in~$s$.  Since $2x+b>0$, by
Lemma~\ref{lem:canon} the substring~$u$ is the standard
representation of~$2x+b$, which we have just shown to be~$w_x b$.
Hence $u=w_x b$, and $s$ contains $w_x b$ as a contiguous
substring.  Let $w_x b=s_i\cdots s_{i+\ell}$.  Deleting the last
character gives $w_x=s_i\cdots s_{i+\ell-1}$ with
$i+\ell-1\le n-1$ (since $i+\ell\le n$), so $w_x$ is a
substring of the prefix~$p$.  This gives $x=\val(w_x)\in V(p)$,
contradicting $x\in V(s)\setminus V(p)$.

\smallskip
Therefore $2x+b\in V(sb)\setminus V(s)$ for each positive
$x\in V(s)\setminus V(p)$.  Distinct values $x\ne x'$ yield
distinct $2x+b\ne2x'+b$ (the map $x\mapsto 2x+b$ is injective on
the integers), so
\begin{equation}\label{eq:append-gain}
  f(sb)-f(s)
  \;\ge\;
  |\{x\in V(s)\setminus V(p):x>0\}|
  \;=\;
  |V(s)\setminus V(p)|-\mathbf{1}[0\in V(s)\setminus V(p)],
\end{equation}
where $\mathbf{1}[\cdot]$ denotes the indicator.

\medskip
\textbf{Step~3:\ the generic case ($0\notin V(s)\setminus V(p)$).}

If $0\notin V(s)\setminus V(p)$, then every element of
$V(s)\setminus V(p)$ is positive, and~\eqref{eq:append-gain}
becomes
\[
  f(sb)-f(s)\ge|V(s)\setminus V(p)|\ge\Delta a(n)
  \qquad\text{(by~\eqref{eq:Vdiff-p})}.
\]
Hence $a(n+1)\ge f(sb)\ge a(n)+\Delta a(n)$, i.e.\
$\Delta a(n+1)\ge\Delta a(n)$.

\medskip
\textbf{Step~4:\ the exceptional case ($0\in V(s)\setminus V(p)$).}

Assume $0\in V(s)\setminus V(p)$.  We show this forces a very
specific structure and then switch to the suffix argument.

Since $0\in V(s)$, the string~$s$ contains at least one~$\mathtt{0}$.
Since $0\notin V(p)$, the prefix $p=s_1\cdots s_{n-1}$ contains
\emph{no} zero, so $p=\mathtt{1}^{n-1}$.  Combined with the existence
of a zero in~$s$ and $s_1=\cdots=s_{n-1}=1$, we must have $s_n=0$, so
\begin{equation}\label{eq:special-s}
  s\;=\;\mathtt{1}^{n-1}\mathtt{0}.
\end{equation}

The suffix is $q=s_2\cdots s_n=\mathtt{1}^{n-2}\mathtt{0}$, which
contains a zero, so $0\in V(q)$.  Therefore
\begin{equation}\label{eq:0-not-in-diff-q}
  0\;\notin\; V(s)\setminus V(q),
\end{equation}
and every element of $V(s)\setminus V(q)$ is positive.

\medskip
\textbf{Step~5:\ prepending a $\mathtt{1}$ creates a new value.}

Let $x\in V(s)\setminus V(q)$.  By~\eqref{eq:0-not-in-diff-q}, $x>0$.
Its canonical witness $w_x$ is the standard representation of~$x$
(Lemma~\ref{lem:canon}).

\textit{Claim:}\ $w_x$ starts at position~$1$ (the first position
of~$s$).

\textit{Proof.}\ Write $w_x=s_i\cdots s_j$.  If $i\ge2$ then
$w_x$ is a substring of $q=s_2\cdots s_n$, giving $x\in V(q)$,
contradicting $x\in V(s)\setminus V(q)$.  Hence $i=1$.

Form the extended string $\mathtt{1}s=\mathtt{1}\,s_1\cdots s_n$ of
length $n+1$.  Since $w_x$ starts at position~$1$ of~$s$, the word
$\mathtt{1}\,w_x$ occurs as a prefix of~$\mathtt{1}s$ (starting at
position~$1$ of~$\mathtt{1}s$).

Let $\ell=|w_x|$ and define $y=\val(\mathtt{1}\,w_x)=2^{\ell}+x$.

\textit{Claim:}\ $\mathtt{1}\,w_x$ is the standard representation
of~$y$.

\textit{Proof.}\ Since $w_x$ is the standard representation of~$x$
with $\ell$ bits, we have $2^{\ell-1}\le x<2^{\ell}$, so
$2^{\ell}\le y=2^{\ell}+x<2^{\ell+1}$.  Hence the standard
representation of~$y$ has exactly $\ell+1$ bits.  The string
$\mathtt{1}\,w_x$ starts with~$\mathtt{1}$, has $\ell+1$ bits, and
has value~$y$; it is therefore the standard representation.

\textit{Claim:}\ $y\notin V(s)$.

\textit{Proof.}\ Suppose $y\in V(s)$.  Take a minimum-length
representative~$u$ of~$y$ in~$s$.  Since $y>0$, by
Lemma~\ref{lem:canon} $u$ is the standard representation of~$y$,
namely $\mathtt{1}\,w_x$.  So $s$ contains $\mathtt{1}\,w_x$ as a
contiguous substring, say $\mathtt{1}\,w_x=s_i\cdots s_{i+\ell}$.
Deleting the first character gives $w_x=s_{i+1}\cdots s_{i+\ell}$
with $i+1\ge2$, so $w_x$ is a substring of $q=s_2\cdots s_n$.  This
gives $x\in V(q)$, contradicting $x\in V(s)\setminus V(q)$.

Therefore $y\in V(\mathtt{1}s)\setminus V(s)$ for each
$x\in V(s)\setminus V(q)$.

\textit{Claim:}\ distinct values of~$x$ yield distinct values of~$y$.

\textit{Proof.}\ Let $x\ne x'$ with canonical witnesses of lengths
$\ell$ and $\ell'$ respectively. Then $y=2^{\ell}+x$ and
$y'=2^{\ell'}+x'$.  If $\ell=\ell'$, then $y\ne y'$ since $x\ne x'$.
If $\ell<\ell'$, then $y=2^{\ell}+x<2^{\ell+1}\le2^{\ell'}\le y'$.
In both cases $y\ne y'$.

\smallskip
Hence
\[
  f(\mathtt{1}s)-f(s)
  \;\ge\;
  |V(s)\setminus V(q)|
  \;\ge\;
  \Delta a(n)
  \qquad\text{(by~\eqref{eq:Vdiff-q})},
\]
and $a(n+1)\ge f(\mathtt{1}s)\ge a(n)+\Delta a(n)$, i.e.\
$\Delta a(n+1)\ge\Delta a(n)$.

\medskip
In both the generic case (Step~3) and the exceptional case (Steps~4--5),
we have shown $\Delta a(n+1)\ge\Delta a(n)$, completing the proof.
\end{proof}

\begin{remark}\label{rem:convex-upper}
Theorem~\ref{thm:convex} shows $\Delta^2 a(n)\ge0$.  Computational
data~\cite{github} further suggest $\Delta^2 a(n)\le2$ for all
$n\ge2$: the second differences take only the values $\{0,1,2\}$ for
all $n\le200$, with the value $2$ occurring at seven values of~$n$
($52, 71, 117, 157, 178, 191, 194$).  Proving the upper bound
$\Delta^2 a(n)\le2$ remains open.
\end{remark}

\section{Computational evidence}\label{sec:comp}

\paragraph{Known exact values ($n\le80$).}
Exact values for $n\le80$ appear in the OEIS b-file submitted by
M.\ Fuller in 2008/9~\cite{oeis}.  No algorithm or proof of completeness
was published; the values have the status of well-verified computational
claims.

\paragraph{Heuristic values ($81\le n\le150$).}
Values of $a(n)$ for $81\le n\le150$ were computed in~\cite{github}
using three independent heuristic methods: unconstrained
Metropolis--Hastings search (at most 10~restarts), constrained
Metropolis--Hastings with a fixed run-length template, and exhaustive
template search over empirically derived block-range constraints.
All three methods agreed on every overlapping case, giving high
confidence in the values; however, these results are \emph{not formally
proved to be exact} and should be understood as heuristic computational
estimates.

\paragraph{Large-scale certified lower bounds.}
A suffix-array-guided local search (\emph{surgical nudge}) was used for
$n$ up to $2\times10^9$~\cite{github}.  In each iteration the algorithm
finds the highest-LCP collision in the suffix array, searches a
neighbourhood for a density-preserving swap that resolves it, and
accepts only swaps that strictly increase the exact score.  Since the
score is monotone non-decreasing, every recorded value is a rigorous
lower bound.  At $n=2\times10^9$ a lower bound of $a(n)/n^2\ge0.49996$
was obtained.

\paragraph{Comparison with theory.}
The ratio $a(n)/n^2$ rises from $0.3596$ at $n=50$ to $0.4031$ at
$n=150$, and exceeds $0.49996$ at $n=2\times10^9$.  
Theorem~\ref{thm:quant} gives $a(n)/n^2\ge1/2-15/\sqrt{n}$, which
reaches $0.4875$ only at $n=14{,}400$.  The computational lower bounds
substantially exceed the theoretical bound at moderate $n$, confirming
that the constant $15$ in $15n^{3/2}$ is far from tight.

\paragraph{Entropy of optimal solutions.}
Let $S_n$ denote the set of all strings achieving $a(n)$ and
$H(n)=\log_2|S_n|$.  For $n\le150$, $H(n)$ is almost always zero or
very small: the typical optimum is unique or nearly so, with
$H(n)\le10.4$ throughout.  Optimal strings share long common prefixes;
the zero-separator structure is completely rigid in 28~cases and varies
by at most one separator in a further 119~cases out of $n\le150$.

\section{Heuristic remarks}\label{sec:heuristic}

To maximise $f(s)$, one seeks to maximise the number of
leading-$\mathtt{1}$ substrings that remain pairwise distinct after
leading zeros are stripped.  Long repeated patterns create large
common-prefix collisions, inflating LCP values and causing many pairs
of substrings to agree.  The explicit construction in
Section~\ref{sec:construction} avoids this by using \emph{strictly
decreasing} one-block lengths: any two substrings $u_{i,j,\alpha,\beta}$
and $u_{i',j',\alpha',\beta'}$ spanning different block-pairs $(i,j)
\ne(i',j')$ have different interior block-length sequences, making them
provably distinct (Lemma~\ref{lem:inject}).

Empirically, all optimal strings for $n\le150$ follow a similar
principle: they begin with a long run of ones, followed by
decreasing-length one-blocks with a fixed separator prefix
$1,2,1,1,\ldots$, and have ones-density growing toward $1$.  This
structure is consistent with the proven density bounds: the empirical
zero count at small $n$ is $O(\sqrt{n})$, well within the theoretical
upper bound of $O(n^{3/4})$ and matching the zero count of the explicit
construction in Section~\ref{sec:construction}.

\section*{Acknowledgements}

The author thanks the contributors to OEIS sequence A112509.

\begin{thebibliography}{9}

\bibitem{oeis}
  OEIS Foundation Inc.,
  \emph{The On-Line Encyclopedia of Integer Sequences},
  Sequence A112509,
  \url{https://oeis.org/A112509}, 2026.

\bibitem{github}
  B.~Clare,
  \emph{A112509: computational exploration of OEIS sequence A112509},
  \url{https://github.com/axelsmagichammer/A112509}, 2026.

\end{thebibliography}

\end{document}
